% =============================================================================
% Bloque 3: Problema (Escasez + Re-muestreo + Deficiencias + Pregunta)
% =============================================================================

\section{Problema}

% --- Slide: Escasez de datos en la práctica (movida desde intro) ---
\begin{frame}{Escasez de datos: un problema cotidiano}
\begin{columns}[T]

\begin{column}{0.5\textwidth}
\small
\textbf{Casos reales que sufren pocos ejemplos:}
\begin{itemize}
  \item Detección de \textbf{discurso de odio} en idiomas minorizados.
  \item Triaje clínico, clasificación de notas médicas.
  \item Categorización de \textbf{tickets de soporte} en empresas.
  \item Spam emergente y nuevos esquemas de fraude.
\end{itemize}
\vspace{0.4em}
\textbf{Por qué duele:}
\begin{itemize}
  \item Etiquetado humano = \alert{caro, lento y sesgado}.
  \item Multiclase $+$ desbalance $\Rightarrow$ macro F1 colapsa.
  \cita{Gopali 2024; Whitehouse 2023}
\end{itemize}
\end{column}

\begin{column}{0.5\textwidth}
\centering
\begin{tikzpicture}[every node/.style={font=\scriptsize}]
  % Iconos representando dominios
  \node[draw=uaiblue, rounded corners, fill=uaibluelight!15, minimum width=2.3cm, minimum height=1cm, text width=2.1cm, align=center] (a) at (-1.3,3) {\textbf{Hate speech}\\\tiny minoría: $<$30 ej.};
  \node[draw=uaiblue, rounded corners, fill=uaibluelight!15, minimum width=2.3cm, minimum height=1cm, text width=2.1cm, align=center] (b) at (1.3,3) {\textbf{Clínico}\\\tiny rare disease};
  \node[draw=uaiblue, rounded corners, fill=uaibluelight!15, minimum width=2.3cm, minimum height=1cm, text width=2.1cm, align=center] (c) at (-1.3,1.5) {\textbf{Soporte TI}\\\tiny producto nuevo};
  \node[draw=uaiblue, rounded corners, fill=uaibluelight!15, minimum width=2.3cm, minimum height=1cm, text width=2.1cm, align=center] (d) at (1.3,1.5) {\textbf{Spam}\\\tiny esquema emergente};
  \node[draw=uaiorange, very thick, rounded corners, fill=uaiorange!15, minimum width=4.8cm, minimum height=0.8cm, text width=4.6cm, align=center] at (0,0.2) {\textbf{Común:} 10--50 ejemplos por clase};
\end{tikzpicture}
\end{column}

\end{columns}
\end{frame}

% --- Slide: Re-muestreo (movida desde Estado del Arte) ---
\begin{frame}{Re-muestreo: el paraguas que une a SMOTE y a los LLMs}
\small
La estadística clásica enmarca la escasez desde hace décadas como un problema de \emph{re-muestreo}. Ambas familias caben aquí.

\vspace{0.5em}
\begin{columns}[T]

\begin{column}{0.50\textwidth}
\setbeamercolor{block title}{fg=white,bg=uaiblue}
\begin{block}{Undersampling}
\scriptsize
\textbf{Reduce} la clase mayoritaria.\\[2pt]
\textbf{Ejemplos:} muestreo aleatorio, NearMiss, Tomek Links.\\[2pt]
\good\ Más rápido al entrenar.\\
\bad\ Pierde información. \alert{Inviable en few-shot}.
\end{block}
\end{column}

\begin{column}{0.50\textwidth}
\setbeamercolor{block title}{fg=white,bg=uaiorange}
\begin{block}{Oversampling}
\scriptsize
\textbf{Aumenta} la clase minoritaria.\\[2pt]
\textbf{Ejemplos:} duplicación, \textbf{SMOTE}, ADASYN, EDA, \textbf{LLMs}.\\[2pt]
\good\ Conserva todos los datos.\\
\bad\ Riesgo de overfitting, ruido o divergencia.
\end{block}
\end{column}

\end{columns}

\vspace{0.7em}
\begin{insight}
\scriptsize
SMOTE y los LLMs caen bajo el mismo paraguas conceptual de \textbf{oversampling}, pero operan en planos distintos: \alert{vectores} vs.\ \alert{texto}.\\[2pt]
La pregunta abierta no es cuál elegir, sino \textbf{cómo combinarlos sin perder garantías geométricas}.
\end{insight}

\mensaje{De ahí surge la pregunta central de esta tesis.}
\end{frame}

% --- Slide: Pregunta de investigación ---
\begin{frame}{Pregunta de investigación}
\vfill
\begin{center}
\setbeamercolor{block title}{fg=white,bg=uaibluedark}
\begin{tcolorbox}[colback=white, colframe=uaibluedark, coltitle=white,
                  title=\centering\textbf{Pregunta central},
                  width=0.9\linewidth, halign=center, arc=4pt, boxrule=2pt]
\Large\itshape ¿Puede un sistema de filtrado geométrico\\
post-generación, desacoplado del LLM,\\[0.4em]
\textbf{mejorar significativamente} la aumentación\\
de datos para clasificación de texto\\
con pocos ejemplos por clase?
\end{tcolorbox}
\end{center}
\vfill
\end{frame}
